\section{Economic Interpretation and Discussion}
\label{sec:discussion}

\subsection{Why forecast performance changes across information stages}
\label{subsec:why_stage_changes}
% Interpret early-stage persistence/factor dominance vs later-stage bridge/data effects.

\subsection{Differences across GDP, inflation, and labour markets}
\label{subsec:cross_domain}
% Relate performance to release frequency, persistence, weekly information,
% target revisions, and contemporaneous indicator availability.

\subsection{Information-set construction versus model sophistication}
\label{subsec:information_vs_model}
% Discuss whether correct real-time design matters more than model complexity.

\subsection{Stable policies versus continual model switching}
\label{subsec:policy_stability}
% Bias-variance / instability / chasing recent performance.

\subsection{Practical implications}
\label{subsec:implications}
% Central banks, government forecasters, financial institutions,
% macro-monitoring systems. Keep claims proportional to evidence.

\subsection{Limitations}
\label{subsec:limitations}
% Candidate limitations:
% - finite real-time evaluation history;
% - imperfect genuine first-release availability for some series;
% - U.S.-only application;
% - pre-specified stage definitions;
% - candidate-model universe;
% - nonlinear alternatives omitted;
% - exceptional COVID observations.
